\bookchapter{第十九章：故障排查与自救}

\chapterintro{手机跟人一样，偶尔也会“生病”。卡顿、闪退、黑屏、充不进电……大部分毛病其实没那么可怕，你完全可以在送去维修之前自己先试试。这章就是你的“手机急诊手册”——先用低成本的方法试，不行再找专业医生。}

\booksection{一、手机变卡、发热——最常被问的问题}

说实话，手机用了一年左右开始变卡，太正常了。原因不一定是手机坏了，更多时候是“积灰”——垃圾文件、后台进程、存储空间慢慢被填满。

\booksubsection{先做诊断，别急着修：}

\begin{enumerate}
  \item \textbf{查存储空间}：去设置 存储看剩余空间。如果不到10\%了，系统会变得非常卡，因为它在挣扎着腾地方。我的经验是：保持至少15\%-20\%的剩余空间，手机才能比较顺畅。
  \item \textbf{查耗电排行}：去电池设置里看过去24小时的耗电排行。如果某个App占了大头但你根本没怎么用，那就是它在后台偷跑。直接卸载或者限制后台活动。
  \item \textbf{联想最近做了什么}：是升级了系统后变卡的？（大版本更新后前几天变卡正常，系统在后台重建索引）。还是安装了某个新App之后变卡的？（卸载试试看）。
  \item \textbf{看使用环境}：手机在太阳底下暴晒、套着厚厚的壳边充边玩游戏，发烫是必然的。高温还会导致CPU自动降频，你就感觉手机突然变卡了。凉下来再试。
\end{enumerate}

\textbf{治疗方法（从简单到复杂）}：重启手机（解决60\%的问题） 清存储空间（删掉没用的App和视频） 某个App搞鬼就清它的缓存或卸载重装 最后一招：备份数据后恢复出厂设置。

\booksection{二、应用闪退、无响应——App不听话了}

某个App点开就闪退、或者用着用着突然卡死不动了。处理流程就像医生问诊，一步步来：

\begin{enumerate}
  \item \textbf{先关掉重开}：从后台划掉App，再重新打开。大多数临时bug这一下就好了。
  \item \textbf{更新App}：去应用商店看有没有更新。有时候是因为App版本太老了，跟系统不兼容。
  \item \textbf{检查权限}：有些App需要某个权限才能正常工作，但你没给。比如相机App没给相机权限打了开就闪退。去设置里找到那个App，把该给的权限都给了。
  \item \textbf{清缓存}：去设置 应用管理 找到那个App 存储 清除缓存。注意是"清除缓存"不是"清除数据"！清除数据会把你App里的登录状态和设置全删了，相当于恢复出厂给App的状态。
  \item \textbf{卸载重装}：最后一招。卸载之前先确认这个App的数据会不会丢——微信的话先备份聊天记录，游戏的话先确认账号有没有绑定手机号。
\end{enumerate}

\booksection{三、充不进电或充电慢——可能不是电池的锅}

充不进电是特别容易让人焦虑的问题，但很多时候原因是出在配件上，不是手机坏了。

\begin{itemize}
  \item \textbf{先排查充电口}：用牙签或者不导电的小工具，轻轻地把充电口里的灰尘和棉絮掏一掏。这个我自己遇到过——口袋里掉了一根线头刚好堵在充电口里，插进去以为是接触不良，掏出来就好了。
  \item \textbf{换根线和充电头试试}：线的损坏率比充电头高得多，尤其是经常卷来卷去的线。换一条确定没问题的线和充电头交叉测试，很快能确定是哪个配件坏了。
  \item \textbf{无线充电的问题}：如果你用无线充电，检查一下手机壳是不是太厚了，或者壳后面有没有金属片（磁吸支架之类）。还有，把手机准确地放到充电板的正中间——歪了就可能充不上。
  \item \textbf{充电慢}：确认充电头功率够不够（iPhone想要快充至少需要20W的PD充电头），线支不支持快充（不是所有的线都能过3A以上电流），手机是不是太烫（太烫会自动降功率保护电池）。
\end{itemize}

\booksection{四、网络连不上、Wi-Fi老是断}

在外面用流量好好的，回家连不上Wi-Fi；或者Wi-Fi连着但就是没网。按这个顺序排查：

\begin{enumerate}
  \item \textbf{开关飞行模式}：下拉控制中心，打开飞行模式等5秒再关掉。这等于让手机重新搜一次网，非常有效。
  \item \textbf{Wi-Fi"忘记"再重连}：去Wi-Fi设置里，点那个网络名旁边的感叹号（安卓是齿轮），选"忽略此网络"或"忘记"，然后重新输入密码连接。
  \item \textbf{重启路由器}：很多人家里的路由器几个月没关过，时间久了会不太稳定。拔掉电源等30秒再插回去，大部分Wi-Fi问题就解决了。
  \item \textbf{检查是不是手机的问题}：别的设备连同一个Wi-Fi正常吗？如果都正常就你的手机不行，那手机可能需要重启或者重置网络设置（设置 通用 传输或还原 还原网络设置，不会删你的数据，但会把所有Wi-Fi密码都清了）。
\end{enumerate}

\booksection{五、屏幕问题：触控不灵、闪屏、留印子}

\begin{itemize}
  \item \textbf{触控不灵敏}：先撕掉钢化膜看看是不是膜的锅——便宜钢化膜用久了确实会影响触控。屏幕上有水或手上有油也会导致触控漂移。都不是的话，重启一下再试。
  \item \textbf{屏幕上有印子/残影（烧屏）}：长时间显示同一个画面（比如一直亮屏刷抖音），会在OLED屏幕上留下"残影"。这不是屏幕坏了，只是暂时的。关屏休息半个小时一般能恢复。
  \item \textbf{屏幕闪烁}：如果是低亮度下闪烁，可能是PWM调光的问题，多见于OLED屏幕。去设置里找"DC调光"或"防频闪"选项，打开后会好很多。不是所有手机都有这个功能。
  \item \textbf{屏幕有一块区域完全黑了/亮了}：这个比较严重，大概率是屏幕物理损坏，需要换屏。非官方维修换的屏幕，色彩和亮度通常不如原装，也影响防水。
\end{itemize}

\booksection{六、无法开机、卡在Logo界面——先别慌}

这种情况最容易让人心跳加速，但冷静，先试试这几个：

\begin{itemize}
  \item \textbf{充电半小时再说}：很多时候只是电彻底耗尽了。插上原装充电器，等至少30分钟，再尝试开机。有些手机在电耗尽后需要充到一定电量才会有反应。
  \item \textbf{强制重启}：iPhone全面屏机型按一下音量加、一下音量减、长按侧边电源键不放，直到屏幕出现苹果Logo（可能要按住10-15秒）。安卓通用方法：长按电源键10秒以上，直到震动或亮屏。
  \item \textbf{卡在启动界面反复重启}：这种情况叫"白苹果"（iPhone）或"卡Logo"（安卓）。如果强制重启不行，可能需要进入恢复模式用电脑修复。iPhone连电脑打开iTunes/Finder，进入恢复模式后选"更新"（不是"恢复"），会重装系统但保留数据。安卓进Fastboot或Recovery模式，各品牌组合键不同，建议搜索你手机型号+"进恢复模式"。
  \item \textbf{以上都不行}：那就只能求助专业维修了。数据如果备份过，就不用太担心。
\end{itemize}

\booksection{七、声音问题：喇叭不响、通话声音小}

\begin{itemize}
  \item \textbf{喇叭没声音}：先确定不是静音模式（iPhone侧边的静音拨片）、勿扰模式、或者音量被按到最低了。听起来很傻，但这种情况占了40\%。然后检查蓝牙是不是连着某个耳机或音箱。
  \item \textbf{通话对方听不到你说话}：检查话筒（麦克风）孔有没有堵住——手机底部的小孔很容易被灰尘堵住。用软毛刷轻轻清理。别用针或牙签硬捅，会捅坏防尘网。
  \item \textbf{声音沙哑/破音}：如果是喇叭进水导致的——不要用吹风机热风吹！用纸巾吸干表面水分，然后放着自然晾干。网上说用声波排水App有一定效果（播放低频声音把水振出去），可以试试。
\end{itemize}

\booksection{八、找官方维修，避开那些坑}

如果上面的方法都试过了还是不行，那就得找专业维修了。但维修市场水很深，说说我自己踩过的坑和经验：

\begin{itemize}
  \item \textbf{找真官方，别找"看起来像官方"的}：去品牌官网（Apple.com.cn、Huawei.com.cn等）找"授权服务商"列表。街边挂个苹果Logo的非授权店，用的可能不是原装配件，修完你的防水就没了一刚。
  \item \textbf{送修前必须做的事}：备份所有数据（别指望修的时候数据不会丢），退出所有重要账号（Apple ID、微信、支付宝），关闭"查找我的手机"。有些维修点会要求你提供锁屏密码，这事儿有隐私风险——你最好自己先退出账号再送过去。
  \item \textbf{换屏幕的坑}：非官方屏幕分"原压屏"（原装玻璃换过）、"后压屏"（非原装玻璃）、"组装屏"（全非原装）。价格差好几倍，但显示效果和耐久度也差很多。如果手机还在保修期，老老实实走官方。
  \item \textbf{换电池的坑}：非原装电池可能虚标容量，而且用了半年就衰减得比原装电池两年的损耗还大。苹果换电池官方价格比非官方贵不了太多，建议走官方。
  \item \textbf{数据恢复的坑}：手机进水/摔坏后，去维修如果对方说"数据恢复要加钱"，小心——有些不良商家会故意夸大损坏程度来多收钱。正规的数据恢复是先报价再操作，而且恢复不成功不收费。
\end{itemize}

\begin{tcbwarn}{万能第一招：重启手机}
不管你遇到什么问题——卡顿、闪退、连不上网、声音异常——先重启一下手机。这句话听起来像是敷衍，但以我自己的经验，至少有六成的小毛病重启之后就好了。手机跟电脑一样，跑久了内存里会积累各种乱七八糟的东西，重启就是清空重来。所以，下次手机出问题，先别急着搜“怎么办”，重启完了再看看。

\end{tcbwarn}

\begin{tcbpractice}{现在就记：这些救命操作要烂熟于心}
\begin{enumerate}
  \item \textbf{强制重启}：全面屏iPhone按音量加 音量减 长按电源键到出现Logo。安卓长按电源键10-15秒直到震动。
  \item \textbf{清理充电口}：找一个小工具（牙签或手机取卡针），轻轻清理充电口的灰尘。
  \item \textbf{Wi-Fi重连}：去Wi-Fi设置里"忘记"你当前的Wi-Fi，重新输入密码连接。
  \item \textbf{开/关飞行模式}：下拉控制中心，打开飞行模式，心里数5秒，关掉，看网络是否恢复。
  \item 把这几个操作的步骤截图保存，或者写在备忘录里。遇到问题的时候不用慌慌张张去搜，直接翻备忘录就行。
\end{enumerate}

\end{tcbpractice}
